\documentclass[12pt]{article}
\usepackage[utf8]{inputenc}
\usepackage{bbm}
\usepackage{geometry}
\geometry{a4paper,scale=0.75}
\linespread{1.5}
\usepackage{graphicx} 
\usepackage{float} 
\usepackage{subcaption} 
\usepackage{enumerate}
\usepackage{amsmath}
\usepackage{array}
\newcolumntype{P}[1]{>{\raggedright\arraybackslash}p{#1}}
\usepackage{booktabs}
\usepackage{multirow}
\usepackage{amsfonts}
\usepackage[english]{babel}
\usepackage{amsthm}
\usepackage{dcolumn}
\usepackage{multicol}
\usepackage{stfloats}
\usepackage{placeins}
\usepackage{lscape}
\usepackage[figuresright]{rotating}
\RequirePackage{pdflscape}
\usepackage[toc,page]{appendix}
\usepackage{geometry}
\usepackage{longtable}
\usepackage{comment}
\usepackage{xcolor}
% \usepackage{amsmath}
\usepackage{amssymb}
\usepackage{empheq}
\usepackage{newtxtext,newtxmath}

% -------- enumerated sub-labels (a), (b), … --
\usepackage{enumitem}
\setlist[enumerate,1]{label=(\alph*),ref=\alph*}
% ---------------------------------------------

\usepackage{hyperref}
\hypersetup{hidelinks,
	colorlinks=true,
	allcolors=black,
	pdfstartview=Fit,
	breaklinks=true}
\usepackage{csquotes}
\usepackage{natbib}
\newtheorem{definition}{Definition}
\newtheorem{theorem}{Theorem}
%\newtheorem{proposition}[theorem]{Proposition}
\newtheorem{proposition}{Proposition}[section]
%\newtheorem{lemma}[theorem]{Lemma}
\newtheorem{lemma}{Lemma}[section]
%\newtheorem{corollary}[theorem]{Corollary}
\newtheorem{corollary}{Corollary}[section]
\newtheorem*{remark}{Remark}
\newtheorem{example}{Example}
\newtheorem{exercise}{Exercise}[section]
\numberwithin{equation}{section}
\newtheorem{assumption}{Assumption}[section] % number within sections

\title{Notes on \cite{alfaroFinanceUncertaintyMultiplier2024}}
\author{Harlly Zhou \\ \texttt{hzhouci@ust.hk}}
\date{\today}

\begin{document}
\maketitle

\paragraph{Motivation} 1. Empirics: Causal effect of second-order shock is correlated with the first-order in economic down-turn; 2. Theory: Current models have difficulty explaining the slow recovery after the Great Recession.

\paragraph{Model Overview} DSGE with heterogeneous firms. 1) Real friction: Investment has fixed cost; 2) Financial friction: Lquidity management by saving cash; 3) Stochasticity: Uncertainty and financing costs.

Results: 1. Amplification effect with financial friction: Double the negative impact of unvertainty shocks on investment; 2. Persistence effect with financial friction: Double the duration of the drop; 3. Propogation effect onto financial variables: Affect both investment and cash (accumulation) and equity payout (drop).

\paragraph{Empirical Results} \textcolor{red}{TO DO!}

\paragraph{Model}
\begin{enumerate}
    \item Technology: Firms use physical capital to produce a signle final good with aggregate productivity $X_t$ and firm-specific productivity $z_{j,t}$:
    \begin{align*}
        y_{j, t} = X_t z_{j,t } k_{j,t}^{\alpha}.
    \end{align*}
    
    Productivity terms follow $AR(1)$ process:
    \begin{align*}
        \log X_{t+1} &= (1-\rho^X) \log \bar X + \rho^X \log X_{t} + \sigma^X_{t} \epsilon^X_{t+1}\\
        \log z_{j,t+1} &= \rho^z \log z_{j,t} + \sigma^z_t \epsilon^z_{j,t+1},
    \end{align*}
    where $\epsilon^X_{t+1}$ and $\epsilon^z_{j,t+1}$ are i.i.d. standard normal shocks and $\sigma^X_t$ and $\sigma^z_t$ are macro and micro uncertainty (time-varying conditional volatilities) of the productivity processes.\footnote{In my paper, the modelling of uncertainty is exactly à la this.} These uncertainties follow a two-state Markov process with high and low states.

    Capital accumulation is standard with non-convex adjustment costs.
    
    \item Cash accumulation: Cash investment and return to holding cash (question mark here).
    
    \item External financing: Exist only when payout is negative. External financing condition is a two-state Markov process with high and low states as a multiplier to the required financing.
    
    \item Firm's problem: Maximize payout minus external financing cost with continuation value.
    \item Households: Have access to equity.
    \item Equilibrium: Goods market, equity market, and money market (money supply is exogenous, equal the demand from firms).
\end{enumerate}

\bibliographystyle{../draft/aer}
\bibliography{../draft/network_uncertainty}

\end{document}
